\section{Setting and two reductions}\label{sec:reduction}

This section states the lattice setting and two conditional theorems. Theorem~\ref{thm:reduction} treats one two-point function. Theorem~\ref{thm:conditional} carries a family of local fields to the Osterwalder--Schrader axioms. Both theorems are proved here, with one exception: Step~9 of the proof of Theorem~\ref{thm:conditional} uses the reconstruction theorem as cited and a description of the reconstructed space that is recalled, not read. Neither theorem proves a property of four-dimensional Yang--Mills theory: each moves the whole difficulty into its hypotheses. Table~\ref{red:status} at the end of the section gives the status of each hypothesis for $\SU(2)$.

\subsection{The lattice theory}\label{red:setting}

Let $\Lambda_k=(\Z/2^k\Z)^4$ be the periodic box of side $2^k$ and let $\mathcal{P}_k$ be its set of plaquettes. A configuration assigns $U_b\in\SU(2)$ to each oriented edge $b$, with $U_{-b}=U_b^{-1}$. For a plaquette $p$ let $U_p$ be the ordered product of the four edge variables around $p$. The Wilson action at coupling $\beta>0$ is
\begin{equation}\label{red:action}
  A_{\beta,k}(U)=\beta\sum_{p\in\mathcal{P}_k}\Bigl(1-\tfrac12\Tr U_p\Bigr),
\end{equation}
and the periodic Wilson state is the probability measure $Z^{-1}e^{-A_{\beta,k}(U)}\prod_b dU_b$, with $dU_b$ the Haar measure. Let $\cL_\beta$ be the set of weak limit points, as $k\to\infty$, of these states, regarded as measures on $\SU(2)^{E(\Z^4)}$, where $E(\Z^4)$ is the set of edges of $\Z^4$. The space $\SU(2)^{E(\Z^4)}$ is compact and metrizable, so $\cL_\beta$ is not empty. For each $\beta$ we choose one $\mu_\beta\in\cL_\beta$ and write $\langle\cdot\rangle_\beta$ for its expectation. Nothing below depends on the choice beyond the hypotheses stated.

Coordinates are $x=(x_0,x_1,x_2,x_3)$, with $x_0=t$ the time. Let $W_4$ be the hypercubic group: the $384$ signed permutation matrices acting on $\R^4$ and on $\Z^4$. For $c\in\Z$ the reflection in the site hyperplane $\{x_\mu=c\}$ acts on configurations, and $\Theta$ denotes the induced map on functions of the configuration. We write $\theta$ for the reflection $(t,\vec x)\mapsto(-t,\vec x)$ in the plane $\{t=0\}$.

\emph{Reflection positivity and the transfer matrix.} Hypothesis~\ref{hyp:H0} below states reflection positivity of $\mu_\beta$ in site hyperplanes. Given it, the standard construction \cite{OS78} produces the following objects; we cite it and do not repeat it, and \cite{OS78} was not read for this section. Let $\mathcal{A}_+$ be the bounded functions of the edge variables of edges with both endpoints in $\{t\ge0\}$. The form $(F,G)\mapsto\langle\overline{\Theta F}\,G\rangle_\beta$ is positive semidefinite on $\mathcal{A}_+$. Let $\cH_\beta$ be the completion of $\mathcal{A}_+$ modulo its null space, $[F]$ the class of $F$, and $\Omega_\beta=[1]$. Let $\tau_n$ be the translation by $n$ lattice steps in the time direction. Then $T_\beta[F]:=[\tau_1F]$ defines an operator with the following properties: $T_\beta$ is self-adjoint by translation invariance; $\langle v,T_\beta^2v\rangle\ge0$ by reflection positivity in the site plane $\{t=1\}$; $\|T_\beta\|\le1$; and $T_\beta\Omega_\beta=\Omega_\beta$. Site-plane reflection positivity alone does not give $T_\beta\ge0$. For the Wilson action $T_\beta\ge0$ is a theorem of \cite{Lu77}, recalled and not read for this section. This paper does not use it.

\begin{definition}\label{red:def-gap}
The \emph{gap in lattice units} is $m(\beta):=-\log\|T_\beta|_{\Omega_\beta^\perp}\|\in[0,\infty]$. When $0<m(\beta)<\infty$, the \emph{correlation length} is $\xi(\beta):=1/m(\beta)$ and the \emph{unit of length} is
\[
  a(\beta):=m(\beta).
\]
A lattice point $n\in\Z^4$ sits at the physical point $a(\beta)n\in a(\beta)\Z^4$.
\end{definition}

\emph{What this choice of unit does.} The unit of length is the correlation length. In this unit the gap equals $1$ by definition, at every $\beta$. The reduction therefore moves nothing. It does not convert a hard statement into an easier one. It fixes the form in which the difficulty must be met, and its content is its hypotheses. Hypothesis~\ref{hyp:A} at weak coupling, $m(\beta)>0$ for large $\beta$ in infinite volume, is the lattice form of the mass gap problem. Hypothesis~\ref{hyp:B} says that the unit $a(\beta)$ shrinks to zero, so that a limit $\beta\to\infty$ is a continuum limit. The other hypotheses ask that correlations stay regular when lengths are measured in this unit.

\subsection{One two-point function}\label{red:R}

Theorem~\ref{thm:reduction} is stated for abstract data, as in the source. For each $\beta\ge\beta_1$ we are given:
\begin{itemize}
\item a Hilbert space $\cH_\beta$ with a unit vector $\Omega_\beta$;
\item a self-adjoint operator $T_\beta$ with $0<T_\beta\le1$ (that is, $T_\beta\ge0$ and $T_\beta$ injective) and $T_\beta\Omega_\beta=\Omega_\beta$; we write $T_\beta=e^{-H_\beta}$ with $H_\beta\ge0$ self-adjoint;
\item a vector $\psi_\beta\ne0$ with $\psi_\beta\perp\Omega_\beta$.
\end{itemize}
The gap is $m(\beta):=\inf\operatorname{spec}(H_\beta|_{\Omega_\beta^\perp})$, which agrees with Definition~\ref{red:def-gap}. Let $\nu_\beta$ be the spectral measure of $\psi_\beta$ for $H_\beta$. It is a finite positive measure on $[m(\beta),\infty)$ with $\nu_\beta([m(\beta),\infty))=\|\psi_\beta\|^2>0$. For real $n\ge0$ put
\begin{equation}\label{red:C}
  C_\beta(n):=\langle\psi_\beta,e^{-nH_\beta}\psi_\beta\rangle=\int e^{-n\varepsilon}\,d\nu_\beta(\varepsilon).
\end{equation}
In the application $\psi_\beta=(O-\langle O\rangle_\beta)\Omega_\beta$ for a bounded gauge-invariant function $O$ of the edges of one time slice near the origin, and $C_\beta(n)$ is its connected correlation at time separation $n$. It is a lattice correlation only at integer $n$; at other $n$ it is defined by the functional calculus. To apply the theorem to the Wilson transfer matrix one needs $0<T_\beta$, which is the theorem of \cite{Lu77}. Theorem~\ref{thm:conditional} does not need it.

\begin{hypothesis}[(A): gap at every coupling]\label{hyp:A}
$m(\beta)>0$ for all $\beta\ge\beta_1$.
\end{hypothesis}

\begin{hypothesis}[(B): criticality at weak coupling]\label{hyp:B}
$m(\beta)\to0$ as $\beta\to\infty$.
\end{hypothesis}

Under (A) and (B), after enlarging $\beta_1$, $\xi(\beta)=1/m(\beta)\in(1,\infty)$ for $\beta\ge\beta_1$.

\begin{hypothesis}[(C): bounded energy at the correlation time]\label{hyp:C}
There are $c\in(0,1)$ and $D<\infty$ such that $C_\beta(c\,\xi)\le D\,C_\beta(\xi)$ for all $\beta\ge\beta_1$.
\end{hypothesis}

The source first stated (C) with $c=\tfrac12$. Every $c\in(0,1)$ gives the same proof, and we state that form. In Theorem~\ref{thm:conditional}, Hypothesis~(A) is called H1 and Hypothesis~(B) is called H2. There $m(\beta)$ is the gap of Definition~\ref{red:def-gap}.

Let $E:=\varepsilon\,\xi$ be the energy in units of the gap. Let $\tilde\nu_\beta$ be the image of $\nu_\beta$ under $\varepsilon\mapsto\varepsilon\xi$; it is carried by $[1,\infty)$. Define the probability measure $\mu_\beta$ on $[1,\infty)$ by
\begin{equation}\label{red:mu}
  d\mu_\beta(E):=\frac{e^{-E}\,d\tilde\nu_\beta(E)}{\int e^{-E'}\,d\tilde\nu_\beta(E')},
  \qquad
  G_\beta(t):=\frac{C_\beta(t\xi)}{C_\beta(\xi)}\quad(t\ge0).
\end{equation}
Hypothesis~(C) says exactly that $\int e^{(1-c)E}\,d\mu_\beta(E)\le D$.

\begin{theorem}[Theorem R$'$]\label{thm:reduction}
Assume Hypotheses~\ref{hyp:A}, \ref{hyp:B} and~\ref{hyp:C}, with constants $c\in(0,1)$ and $D$. Then:
\begin{enumerate}
\item[(R1)] For every $\beta\ge\beta_1$, $\mu_\beta$ is a probability measure on $[1,\infty)$ and $G_\beta(t)=\int e^{-(t-1)E}\,d\mu_\beta(E)$ for every $t\ge0$.
\item[(R2)] $\mu_\beta((K,\infty))\le D\,e^{-(1-c)K}$ for every $K\ge0$, and $\int E\,d\mu_\beta\le \log D/(1-c)$, for every $\beta\ge\beta_1$.
\item[(R3)] Every sequence $\beta_k\to\infty$ has a subsequence along which $\mu_{\beta}$ converges weakly to a probability measure $\mu$ on $[1,\infty)$, and $G_\beta(t)\to G(t):=\int e^{-(t-1)E}\,d\mu(E)$ for every $t>c$.
\item[(R4)] $G(1)=1$; $0<G(t)\le e^{-(t-1)}$ for $t\ge1$; $\int E\,d\mu\le\log D/(1-c)$; and $G(t)\ge D^{-(t-1)/(1-c)}$ for $t\ge1$.
\item[(R5)] The unit of length $a(\beta)=m(\beta)$ tends to $0$.
\end{enumerate}
\end{theorem}

\begin{proof}
(R1). By (A), $\xi<\infty$ and $\tilde\nu_\beta$ is carried by $[1,\infty)$. The denominator in \eqref{red:mu} equals $C_\beta(\xi)=\int e^{-\xi\varepsilon}\,d\nu_\beta>0$, because $\nu_\beta\ne0$ and the integrand is positive. So $\mu_\beta$ is a probability measure on the closed set $[1,\infty)$. For $t\ge0$,
\[
  \frac{C_\beta(t\xi)}{C_\beta(\xi)}=\frac{\int e^{-tE}\,d\tilde\nu_\beta}{\int e^{-E}\,d\tilde\nu_\beta}=\int e^{-(t-1)E}\,d\mu_\beta(E).
\]
Both integrals are finite at fixed $\beta$ because $\tilde\nu_\beta$ is finite and $e^{-tE}\le1$. Hypothesis~(C) is used only for bounds uniform in $\beta$.

(R2). For $E>K$ we have $e^{-E}\le e^{-(1-c)K}e^{-cE}$. Hence
\[
  \int_{(K,\infty)}e^{-E}\,d\tilde\nu_\beta\le e^{-(1-c)K}\int e^{-cE}\,d\tilde\nu_\beta=e^{-(1-c)K}C_\beta(c\xi)\le D\,e^{-(1-c)K}C_\beta(\xi).
\]
Divide by $C_\beta(\xi)$. For the mean, Jensen's inequality for the convex function $\exp$ gives $\exp\bigl((1-c)\int E\,d\mu_\beta\bigr)\le\int e^{(1-c)E}\,d\mu_\beta=C_\beta(c\xi)/C_\beta(\xi)\le D$.

(R3). By (R2) the family $(\mu_\beta)$ is tight on $[1,\infty)$. Prokhorov's theorem gives a weakly convergent subsequence. Its limit $\mu$ has mass $1$, and it is carried by $[1,\infty)$ because that set is closed. For $t\ge1$ the function $E\mapsto e^{-(t-1)E}$ is bounded and continuous on $[1,\infty)$, so the integrals converge. For $c<t<1$ the integrand $e^{(1-t)E}$ is continuous and uniformly integrable:
\[
  \int_{(K,\infty)}e^{(1-t)E}\,d\mu_\beta\le e^{-(t-c)K}\int e^{(1-c)E}\,d\mu_\beta\le D\,e^{-(t-c)K}.
\]
Truncation at $K$ and weak convergence then give $G_\beta(t)\to G(t)$.

(R4). $G(1)=\mu([1,\infty))=1$. For $t\ge1$ and $E\ge1$ we have $0<e^{-(t-1)E}\le e^{-(t-1)}$, and the integral of a positive function against a probability measure is positive. The function $E\mapsto E$ is continuous and bounded below, so $\int E\,d\mu\le\liminf\int E\,d\mu_\beta\le\log D/(1-c)$ along the subsequence. Jensen's inequality gives $G(t)\ge\exp\bigl(-(t-1)\int E\,d\mu\bigr)\ge D^{-(t-1)/(1-c)}$ for $t\ge1$.

(R5). This is Definition~\ref{red:def-gap} and Hypothesis~\ref{hyp:B}.
\end{proof}

Hypothesis~(B) enters only in (R5). Theorem~\ref{thm:reduction} is Prokhorov's theorem and properties of Laplace transforms. The conclusions ``gap at least $1$'' and ``not zero'' hold by construction: the unit is the gap, and the correlation is normalized at one correlation length.

Every admissible constant in Hypothesis~\ref{hyp:C} satisfies $D\ge e^{1-c}$, so $D\ge\sqrt e$ when $c=\tfrac12$. Indeed $\mu_\beta$ is a probability measure carried by $E\ge1$, so $D\ge\int e^{(1-c)E}\,d\mu_\beta\ge e^{1-c}$.

\emph{What Theorem~\ref{thm:reduction} does not give.} It gives one two-point function of one observable along one time axis, along subsequences. It gives no Euclidean invariance and no interaction: a free massive lattice field satisfies its hypotheses, and the theorem returns a free field. The lowest energy that $\psi_\beta$ sees in the limit can exceed $1$, so the theorem shows a gap of at least $1$, not equal to $1$. The limit can be $\mu=\delta_1$, a pure exponential, which a local relativistic field does not produce at a fixed spatial point. The normalization $C_\beta(\xi)$ tends to zero, and nothing in the theorem makes the normalizations of different observables consistent.

Hypothesis~\ref{hyp:C} is not a statement about locality alone. It also asks that the observable overlap the lightest states, as the following example shows.

\begin{observation}\label{red:twomass}
Let the data be two decoupled free massive lattice scalar fields on $\Z^4$ with masses $m_1$ fixed and $m_2=m_2(\beta)\to0$, and let $\psi_\beta=\varphi_1(0)\Omega_\beta$, the first field at one site. Then Hypotheses~\ref{hyp:A} and~\ref{hyp:B} hold, and Hypothesis~\ref{hyp:C} fails for every $c\in(0,1)$ and every $D<\infty$.
\end{observation}

\begin{proof}
For one free field of mass $m$ the one-particle energies are $2\operatorname{asinh}\bigl(\tfrac12(m^2+\sum_{i=1}^3 4\sin^2(p_i/2))^{1/2}\bigr)\ge E_0(m):=2\operatorname{asinh}(m/2)$, and $\varphi(0)\Omega$ lies in the one-particle space. For the decoupled pair the gap is $m(\beta)=\min(E_0(m_1),E_0(m_2))=E_0(m_2)$, which is positive and tends to $0$. The spectral measure of $\psi_\beta$ is carried by $[E_0(m_1),\infty)$. Hence
\[
  \frac{C_\beta(c\xi)}{C_\beta(\xi)}=\frac{\int e^{-c\xi\varepsilon}\,d\nu_\beta}{\int e^{-\xi\varepsilon}\,d\nu_\beta}\ge e^{(1-c)\xi E_0(m_1)}\to\infty
\]
as $\xi=1/E_0(m_2)\to\infty$.
\end{proof}

The free-field dispersion relation used here is standard; Section~\ref{red:free} records its numerical check.

\subsection{Fields and hypotheses for the full chain}\label{red:S-setting}

Fix a finite family $O_1,\dots,O_r$ of \emph{site-scalar} observables. Each $O_i(n)$ is a bounded, real, gauge-invariant function of the edges with both endpoints within lattice distance $1$ of the site $n$, and it is invariant under $W_4$ acting about $n$. An example is $O(n)=\sum_p\tfrac12\Tr U_p$ over the $24$ plaquettes having $n$ as a corner. Write $\hat O_i(n):=O_i(n)-\langle O_i\rangle_\beta$. A reflection $\theta'$ in any site hyperplane maps $\hat O_i(n)$ to $\hat O_i(\theta'n)$ exactly, with no shift of the base point. That is the reason for site-scalar fields.

Let $N_\beta:=2\lfloor\xi/2\rfloor$ and define the normalization
\[
  Z_i(\beta)^2:=\langle\hat O_i(0)\,\hat O_i(N_\beta e_0)\rangle_\beta .
\]
For $x=(x_1,\dots,x_k)\in(a\Z^4)^k$ and indices $\iota=(i_1,\dots,i_k)$, the \emph{lattice Schwinger functions} are
\begin{equation}\label{red:Sbeta}
  S^\beta_\iota(x):=\prod_{j=1}^k Z_{i_j}(\beta)^{-1}\,\Bigl\langle\prod_{j=1}^k\hat O_{i_j}(x_j/a)\Bigr\rangle_\beta .
\end{equation}
Let ${}^0\mathcal S(\R^{4k})$ be the space of Schwartz functions that vanish, with all their derivatives, wherever $x_i=x_j$ for some $i\ne j$, and let ${}^0\mathcal S_+(\R^{4k})$ be the subspace of functions whose support lies in $\{0<t_1<\dots<t_k\}$ \cite[eq.~(2.25)]{KQR21}. For $f\in{}^0\mathcal S(\R^{4k})$ put
\[
  S^\beta_\iota(f):=a^{4k}\sum_{x\in(a\Z^4)^k}S^\beta_\iota(x)\,f(x).
\]
A configuration with two coinciding points contributes $0$, because $f$ vanishes there. We use the Schwartz norms $|f|_N:=\max_{|\alpha|\le N}\sup_y(1+|y|)^N|\partial^\alpha f(y)|$. The norms in \cite[eq.~(9.18)]{KQR21} were not compared with these; E0$'$ below is stated with ours. For $x\in(\R^4)^k$ let $r(x):=\min_{i<j}|x_i-x_j|$ and $\rho_j(x):=\min_{l\ne j}|x_j-x_l|$.

There is no smoothing flow. Every point carries a local composite field, renormalized multiplicatively by $Z_i(\beta)$ and evaluated only at non-coinciding points, which is what ${}^0\mathcal S$ allows. All short-distance control is placed in Hypothesis~\ref{hyp:H3}.

\begin{hypothesis}[H0: lattice axioms]\label{hyp:H0}
$\mu_\beta$ is invariant under the translations $\Z^4$ and under $W_4$. It is reflection positive for reflections in every site hyperplane $\{x_\mu=c\}$, $c\in\Z$: $\langle\overline{\Theta F}\,F\rangle_\beta\ge0$ for every bounded $F$ depending only on edges in $\{x_\mu\ge c\}$.
\end{hypothesis}

H1 is Hypothesis~\ref{hyp:A} and H2 is Hypothesis~\ref{hyp:B}, both for the gap $m(\beta)$ of Definition~\ref{red:def-gap} of the state $\mu_\beta$.

\begin{hypothesis}[H3: uniform regularity in units of the gap]\label{hyp:H3}
There are constants $\alpha,\gamma,b,D<\infty$, and for each $k$ and $\delta>0$ a constant $L_{k,\delta}$, such that for all $\beta\ge\beta_1$, all $k$ and all $\iota$:
\begin{enumerate}
\item[(a)] $Z_i(\beta)>0$;
\item[(b)] (growth) $|S^\beta_\iota(x)|\le\alpha\gamma^k(k!)^b\prod_{j=1}^k\max(1,\rho_j(x)^{-D})$ for $x\in(a\Z^4)^k$ with distinct points;
\item[(c)] (equicontinuity) $|S^\beta_\iota(x)-S^\beta_\iota(y)|\le L_{k,\delta}|x-y|$ for $x,y\in(a\Z^4)^k$ with all pairwise distances at least $\delta$.
\end{enumerate}
\end{hypothesis}

Increasing $D$ weakens (b), so we take $D\ge1$.

\begin{hypothesis}[H4: asymptotic rotation invariance, one rotation]\label{hyp:H4}
For one angle $\vartheta$ with $\vartheta/\pi$ irrational, let $R_\vartheta$ be the rotation by $\vartheta$ in the $(x_1,x_2)$-plane, and let $[y]_a$ be the nearest point of $a\Z^4$ to $y$ (ties broken lexicographically). For every $k$, $\iota$, $\delta>0$ and $R<\infty$,
\[
  \sup\bigl|S^\beta_\iota([R_\vartheta x]_a)-S^\beta_\iota([x]_a)\bigr|\to0\qquad(\beta\to\infty),
\]
the supremum taken over configurations in the ball of radius $R$ with pairwise distances at least $\delta$.
\end{hypothesis}

Here $[x]_a$ is taken pointwise. Each point moves by at most $a$, so $|[x]_a-x|\le a\sqrt k$ in $(\R^4)^k$.

\emph{The Osterwalder--Schrader axioms.} We read the axioms in \cite{KQR21}, Section~2.2 and Section~9.1--9.2. We did not read \cite{OS73} or \cite{OS75}. For $f$ on $\R^{4k}$ write $(\Theta f)^*(x):=\overline{f(\theta x_1,\dots,\theta x_k)}$, and for $s\in\R$ write $f_s(x):=f(x_1-se_0,\dots,x_k-se_0)$. A family of distributions $S_\iota$ on ${}^0\mathcal S(\R^{4k})$ satisfies:
\begin{itemize}
\item[E0$'$] (linear growth) if there is an integer $s$ with $|S_\iota(f)|\le\sigma_k|f|_{ks}$ for all $f\in{}^0\mathcal S(\R^{4k})$, where $\sigma_k\le\alpha'(k!)^{b'}$ for constants $\alpha',b'$ \cite[eq.~(9.18)]{KQR21};
\item[E1] (Euclidean covariance) if $S_\iota(f\circ g)=S_\iota(f)$ for every Euclidean motion $g$ of $\R^4$, acting on each point;
\item[E2] (reflection positivity) if $\sum_{\iota,\iota'}S_{(\iota',\iota)}\bigl((\Theta f_{\iota'})^*\otimes f_\iota\bigr)\ge0$ for every finite family $(f_\iota)$ with $f_\iota\in{}^0\mathcal S_+$, where the empty index tuple carries a constant and $S_\emptyset=1$;
\item[E3] (symmetry) if $S_\iota$ is invariant under simultaneous permutation of the points and the indices;
\item[E4] (cluster) if $S_{(\iota',\iota)}\bigl((\Theta f)^*\otimes g_{\lambda\vec v}\bigr)\to S_{\iota'}((\Theta f)^*)\,S_\iota(g)$ as $\lambda\to\infty$, for $f,g\in{}^0\mathcal S_+$ and every spatial vector $\vec v\ne0$; here $g_{\lambda\vec v}$ is $g$ translated by $\lambda\vec v$.
\end{itemize}
In \cite{KQR21}, E4 is named as the Osterwalder--Schrader cluster property for spatial translations (footnote~20 there); the formula displayed is the standard form, written by us. \cite[Section~9.2.1]{KQR21} records that Lemma~8.8 of \cite{OS73} is wrong; the 1975 form with E0$'$ is the one used here.

\begin{theorem}[Reconstruction theorem \cite{OS75}, as read in \cite{KQR21}]\label{red:OS}
Distributions $S_\iota$ satisfying E0$'$ and E1--E3 are the Schwinger functions of a Wightman theory with fields $\varphi_1,\dots,\varphi_r$.
\end{theorem}

We also use the form of the reconstructed objects: the Hilbert space is the completion of finite families in ${}^0\mathcal S_+$ modulo the null space of the E2 form, the vacuum $\Omega$ is the class of the constant $1$, and $e^{-sH}$ acts by $[f]\mapsto[f_s]$ for $s\ge0$. This description is taken from \cite[Section~4]{OS73} as recalled. It was not read.

\begin{theorem}[Theorem S]\label{thm:conditional}
Assume Hypotheses~\ref{hyp:H0}, \ref{hyp:A} (H1), \ref{hyp:B} (H2), \ref{hyp:H3} and~\ref{hyp:H4}. Then there is a sequence $\beta_j\to\infty$ such that, for every $k$ and $\iota$:
\begin{enumerate}
\item $S^{\beta_j}_\iota\to S_\iota$ uniformly on compact sets of configurations with pairwise distances at least $\delta>0$ (through the lattice approximants of Step~1 below), with $S_\iota$ continuous on the set of configurations of pairwise distinct points; and $S^{\beta_j}_\iota(f)\to S_\iota(f):=\int S_\iota f$ for every $f\in{}^0\mathcal S(\R^{4k})$.
\item The $S_\iota$ satisfy the Osterwalder--Schrader axioms in the 1975 form: E0$'$, E1, E2, E3 and E4.
\item $S_{ii}(0,e_0)=1$ for every $i$, so the theory is not empty.
\item By Theorem~\ref{red:OS} there is a Wightman theory with fields $\varphi_1,\dots,\varphi_r$. Its vacuum is the unique translation-invariant vector, and the joint spectrum of $(H,\vec P)$ on $\Omega^\perp$ lies in $\{p^2\ge1,\ p^0>0\}$: the mass gap is at least $1$ in these units.
\end{enumerate}
The theorem asserts nothing about uniqueness of the limit, interaction, short-distance behaviour, or other gauge groups.
\end{theorem}

\subsection{Proof of Theorem~\ref{thm:conditional}}\label{red:S-proof}

By H2, after enlarging $\beta_1$, $a(\beta)\le1$ for $\beta\ge\beta_1$; by H1, $a(\beta)>0$. Fix any sequence $\beta\to\infty$; all subsequences below are taken from it, and $a\to0$ along it.

\emph{Step 1: compactness.} For $x\in(\R^4)^k$ whose rounded points $[x_j]_a$ are distinct, put $\tilde S^\beta_\iota(x):=S^\beta_\iota([x]_a)$. For $\ell\in\N$ let
\[
  K_\ell:=\{x\in(\R^4)^k:\ |x_j|\le\ell\text{ for all }j,\ r(x)\ge1/\ell\}.
\]
If $a\le1/(4\ell)$, then for $x\in K_\ell$ the rounded configuration has pairwise distances at least $1/\ell-2a\ge1/(2\ell)$. By H3(b), $|\tilde S^\beta_\iota|\le\alpha\gamma^k(k!)^b(2\ell)^{Dk}$ on $K_\ell$. By H3(c), for $x,y\in K_\ell$,
\begin{equation}\label{red:asym-eq}
  |\tilde S^\beta_\iota(x)-\tilde S^\beta_\iota(y)|\le L_{k,1/(2\ell)}\bigl(|x-y|+2a\sqrt k\bigr).
\end{equation}

\begin{lemma}\label{red:AA}
Let $K$ be a compact metric space and $f_j:K\to\R$ functions with $\sup|f_j|\le M$ and $|f_j(x)-f_j(y)|\le L\,d(x,y)+\varepsilon_j$, where $\varepsilon_j\to0$. Then a subsequence converges uniformly on $K$ to an $L$-Lipschitz function.
\end{lemma}

\begin{proof}
Let $Q\subset K$ be countable and dense. A diagonal argument gives a subsequence, still written $f_j$, that converges at every point of $Q$. Given $\eta>0$, choose a finite set $Q_\eta\subset Q$ such that every $x\in K$ lies within $\eta$ of some $q\in Q_\eta$. Then $|f_i(x)-f_j(x)|\le2L\eta+\varepsilon_i+\varepsilon_j+\max_{q\in Q_\eta}|f_i(q)-f_j(q)|$, which is less than $(2L+1)\eta$ for $i,j$ large. So the sequence is uniformly Cauchy. Letting $j\to\infty$ in the hypothesis shows that the limit is $L$-Lipschitz.
\end{proof}

Apply Lemma~\ref{red:AA} on each $K_\ell$ with $\varepsilon=2aL_{k,1/(2\ell)}\sqrt k$, and extract diagonally over the countably many triples $(k,\iota,\ell)$. This gives $\beta_j\to\infty$ and functions $S_\iota$ on the set of configurations of distinct points, Lipschitz on each $K_\ell$, with $\tilde S^{\beta_j}_\iota\to S_\iota$ uniformly on each $K_\ell$. Every compact set of configurations with pairwise distances at least $\delta$ lies in some $K_\ell$. Since $\rho_j([x]_a)\ge\rho_j(x)-2a$, H3(b) passes to the limit:
\begin{equation}\label{red:growth-limit}
  |S_\iota(x)|\le\alpha\gamma^k(k!)^b\prod_{j}\max(1,\rho_j(x)^{-D})\quad\text{for distinct points.}
\end{equation}

\emph{Step 2: the distribution property and E0$'$.}

\begin{lemma}[Taylor]\label{red:taylor}
Let $f\in{}^0\mathcal S(\R^{4k})$, $M\ge1$ and $P\ge0$, and let $N\ge\max(M,P)$. If $r(x)\le1$, then
\[
  |f(x)|\le\frac{(4k)^M}{M!}\,2^P\,|f|_N\;r(x)^M\,(1+|x|)^{-P}.
\]
\end{lemma}

\begin{proof}
Choose $i<j$ with $|x_i-x_j|=r(x)$. Let $z$ agree with $x$ except that $z_i=z_j=\tfrac12(x_i+x_j)$. Then $z$ lies in $\{x_i=x_j\}$ and $|x-z|=r(x)/\sqrt2<1$. Put $g(s):=f(z+s(x-z))$ for $s\in[0,1]$. All derivatives of $f$ of order less than $M$ vanish at $z$, so Taylor's formula with integral remainder gives
\[
  f(x)=g(1)=\int_0^1\frac{(1-s)^{M-1}}{(M-1)!}\,g^{(M)}(s)\,ds,
  \qquad
  g^{(M)}(s)=\sum_{|\alpha|=M}\frac{M!}{\alpha!}\,\partial^\alpha f(z+s(x-z))\,(x-z)^\alpha .
\]
Now $|(x-z)^\alpha|\le|x-z|^M\le r(x)^M$ and $\sum_{|\alpha|=M}M!/\alpha!=(4k)^M$. Every point $y$ of the segment satisfies $|y-x|<1$, so $1+|x|\le2(1+|y|)$ and $|\partial^\alpha f(y)|\le|f|_N(1+|y|)^{-P}\le2^P|f|_N(1+|x|)^{-P}$. Combining, $|f(x)|\le\frac{1}{M!}(4k)^M r(x)^M2^P|f|_N(1+|x|)^{-P}$.
\end{proof}

Take $M:=Dk+1$ and $P:=4k+2$. Then $N:=(D+6)k\ge\max(M,P)$. Put
\[
  A_k:=\alpha\gamma^k(k!)^b\,2^{4k+2}\,(4e)^{Dk+1}.
\]
Since $M!\ge(M/e)^M$ and $D\ge1$, we have $(4k)^M/M!\le(4ek/M)^M\le(4e)^{M}$.

\emph{Claim.} Let $x$ have distinct points, and let $S$ be either $S^\beta_\iota$ (for $x$ in the lattice) or $S_\iota$. Then
\begin{equation}\label{red:pointwise}
  |S(x)f(x)|\le A_k\,|f|_{(D+6)k}\,\min(1,r(x))\,(1+|x|)^{-(4k+2)} .
\end{equation}
If $r(x)\ge1$, every $\rho_j\ge1$, so $|S(x)|\le\alpha\gamma^k(k!)^b$ by H3(b) or \eqref{red:growth-limit}, and $|f(x)|\le|f|_N(1+|x|)^{-P}$. If $r(x)<1$, then $\prod_j\max(1,\rho_j^{-D})\le r(x)^{-Dk}$, and Lemma~\ref{red:taylor} gives $|f(x)|\le(4e)^M2^P|f|_Nr(x)^{Dk+1}(1+|x|)^{-P}$. The product is at most the right side of \eqref{red:pointwise}.

Let $I_k:=\int_{\R^{4k}}(1+|y|)^{-(4k+2)}\,dy=|S^{4k-1}|/\bigl(4k(4k+1)\bigr)$ and $I'_k:=\int_{\R^{4k}}(1+|y|)^{-(4k+1)}\,dy$; both are finite. If $a\sqrt{4k}\le1$, then each lattice point $x$ is within distance $1$ of every point $y$ of the cube $x+[0,a)^{4k}$, so $1+|y|\le2(1+|x|)$ and
\begin{equation}\label{red:riemann}
  a^{4k}\sum_{x\in(a\Z^4)^k}(1+|x|)^{-P}\le2^P\int_{\R^{4k}}(1+|y|)^{-P}\,dy .
\end{equation}

\emph{E0$'$.} Integrating \eqref{red:pointwise} gives $|S_\iota(f)|\le A_kI_k|f|_{(D+6)k}$. With $|S^{4k-1}|=2\pi^{2k}/\Gamma(2k)\le3\pi^{2k}$ and $4k(4k+1)\ge20$,
\[
  A_kI_k\le7\alpha\,\Gamma^k(k!)^b,\qquad \Gamma:=16\gamma(4e)^D\pi^2,
\]
and $\Gamma^k\le e^\Gamma k!$. This is E0$'$ with $s=D+6$ and $\sigma_k\le7\alpha e^{\Gamma}(k!)^{b+1}$. For the lattice functions, \eqref{red:pointwise} and \eqref{red:riemann} give
\begin{equation}\label{red:uniform}
  |S^\beta_\iota(f)|\le2^{4k+2}A_kI_k\,|f|_{(D+6)k}\qquad\text{whenever }a\le1/(2\sqrt k).
\end{equation}

\emph{Convergence of $S^{\beta_j}_\iota(f)$.} Fix $f$, $\delta\in(0,1)$ and $R\ge1$. Let $\mathsf A:=\{x:|x|\le R,\ r(x)\ge\delta\}$ and let $\mathsf B$ be its complement. On $\mathsf B$, either $r(x)<\delta$, or $|x|>R$ and then $(1+|x|)^{-(4k+2)}\le(1+R)^{-1}(1+|x|)^{-(4k+1)}$. By \eqref{red:pointwise} and \eqref{red:riemann}, the contributions of $\mathsf B$ to $S^{\beta_j}_\iota(f)$ and to $\int S_\iota f$ are each at most $2^{4k+2}A_k|f|_{(D+6)k}(\delta I_k+(1+R)^{-1}I'_k)$, uniformly in $j$ large. The set $\mathsf A$ is compact and lies in some $K_\ell$, on which $\tilde S^{\beta_j}_\iota\to S_\iota$ uniformly; at lattice points $\tilde S^{\beta_j}_\iota=S^{\beta_j}_\iota$. So $a^{4k}\sum_{x\in\mathsf A}|S^{\beta_j}_\iota-S_\iota||f|\to0$. Finally $S_\iota f$ is continuous on a neighbourhood of $\mathsf A$, and the boundary of $\mathsf A$ has Lebesgue measure zero, so the Riemann sums $a^{4k}\sum_{x\in\mathsf A}S_\iota(x)f(x)$ converge to $\int_{\mathsf A}S_\iota f$. Letting $j\to\infty$, then $\delta\to0$ and $R\to\infty$, gives $S^{\beta_j}_\iota(f)\to S_\iota(f)$. This proves part~1 of the theorem and E0$'$.

\emph{Step 3: translations.} By H0, $S^\beta_\iota(x+an)=S^\beta_\iota(x)$ for $n\in\Z^4$, and $[x+an]_a=[x]_a+an$. Let $y\in\R^4$ and choose $n_j\in\Z^4$ with $|a_jn_j-y|\le a_j$. By \eqref{red:asym-eq}, $|\tilde S^{\beta_j}_\iota(x+y)-\tilde S^{\beta_j}_\iota(x+a_jn_j)|\le L(a_j\sqrt k+2a_j\sqrt k)\to0$ on each $K_\ell$, and $\tilde S^{\beta_j}_\iota(x+a_jn_j)=\tilde S^{\beta_j}_\iota(x)$. Hence $S_\iota(x+y)=S_\iota(x)$.

\emph{Step 4: rotations.} Let $\mathcal G$ be the set of $g\in O(4)$ with $S_\iota(gx)=S_\iota(x)$ for all $k$, $\iota$ and all configurations $x$ of distinct points. It is a subgroup, and it is closed because each $S_\iota$ is continuous on the $O(4)$-invariant open set of distinct configurations. $W_4\subset\mathcal G$: by H0, $\mu_\beta$ is $W_4$-invariant, $a\Z^4$ is $W_4$-invariant and each $O_i$ is a $W_4$-scalar, so $S^\beta_\iota(wx)=S^\beta_\iota(x)$ for lattice $x$ and $w\in W_4$; the rounding $[wx]_a$ differs from $w[x]_a$ by at most $2a\sqrt k$, which costs $O(a)$ by \eqref{red:asym-eq}. $R_\vartheta\in\mathcal G$: $R_\vartheta$ is an isometry, so it maps the ball of radius $R$ with separation $\delta$ to itself, and H4 together with the uniform convergence of Step~1 gives $S_\iota(R_\vartheta x)=\lim\tilde S^{\beta_j}_\iota(R_\vartheta x)=\lim\tilde S^{\beta_j}_\iota(x)=S_\iota(x)$. By Lemma~\ref{red:dense} below, $\mathcal G=O(4)$. With Step~3 and the change of variables $\int S_\iota(x)f(g^{-1}x)\,dx=\int S_\iota(gx)f(x)\,dx$, this is E1.

\begin{lemma}\label{red:dense}
Let $\vartheta/\pi$ be irrational. The closed subgroup $\mathcal K$ of $O(4)$ generated by $W_4$ and $R_\vartheta$ is $O(4)$.
\end{lemma}

\begin{proof}
For $0\le\mu<\nu\le3$ let $L_{\mu\nu}:=E_{\mu\nu}-E_{\nu\mu}$, where $E_{\mu\nu}$ is the matrix unit, and let $R^{\mu\nu}_\varphi:=\exp(\varphi L_{\mu\nu})$, the rotation by $\varphi$ in the $(x_\mu,x_\nu)$-plane. So $R_\vartheta=R^{12}_\vartheta$.

(i) Kronecker. Since $\vartheta/(2\pi)$ is irrational, the multiples $n\vartheta$, $n\in\N$, are dense modulo $2\pi$. So the powers $R_\vartheta^n=R^{12}_{n\vartheta}$ are dense in $\{R^{12}_\varphi:\varphi\in\R\}$, and $\mathcal K$, being closed, contains every $R^{12}_\varphi$.

(ii) Conjugation. For $\mu<\nu$ let $\sigma\in W_4$ be a permutation matrix with $\sigma e_1=e_\mu$ and $\sigma e_2=e_\nu$. Then $\sigma R^{12}_\varphi\sigma^{-1}=R^{\mu\nu}_\varphi$, so $\mathcal K$ contains the six one-parameter groups $\{R^{\mu\nu}_\varphi\}$.

(iii) The six matrices $L_{\mu\nu}$ form a basis of the antisymmetric $4\times4$ matrices, which is the Lie algebra $\mathfrak{so}(4)$, of dimension $6$. Let $\Phi(\varphi_1,\dots,\varphi_6)$ be the product of the six $\exp(\varphi_pL_p)$ in a fixed order. Then $\Phi:\R^6\to SO(4)$ is smooth with $\Phi(0)=1$, and its differential at $0$ is $(\varphi_p)\mapsto\sum_p\varphi_pL_p$, an isomorphism onto $\mathfrak{so}(4)$. By the inverse function theorem $\Phi$ maps a neighbourhood of $0$ onto a neighbourhood $U$ of $1$ in $SO(4)$, and $U\subset\mathcal K$ by (ii).

(iv) Connectedness. The subgroup generated by $U$ is a union of open sets, hence open; an open subgroup is also closed, since its complement is a union of cosets. $SO(4)$ is connected, so this subgroup is $SO(4)$, and $SO(4)\subset\mathcal K$.

(v) Reflection. $W_4$ contains $\operatorname{diag}(-1,1,1,1)$, of determinant $-1$. So $\mathcal K\supset SO(4)\cup\operatorname{diag}(-1,1,1,1)SO(4)=O(4)$.
\end{proof}

\emph{Step 5: symmetry (E3) and reality.} The lattice fields are commuting real random variables. So $S^\beta_\iota(x)$ is real and invariant under simultaneous permutation of points and indices. Both properties pass to pointwise limits.

\emph{Step 6: reflection positivity (E2).} Let $\mathcal T(\R^{4k})$ be the set of $f\in{}^0\mathcal S(\R^{4k})$ whose (closed) support lies in $\{x: t(x_j)>0\text{ for all }j\}$. Then ${}^0\mathcal S_+\subset\mathcal T$. We prove positivity on $\mathcal T$, which we need in Step~7. Let $f=(f_\iota)$ be a finite family with $f_\iota\in\mathcal T$ (and a constant for the empty tuple), and set
\[
  F_\beta:=\sum_\iota a^{4k}\sum_{x\in(a\Z^4)^k}f_\iota(x)\prod_{j}Z_{i_j}(\beta)^{-1}\hat O_{i_j}(x_j/a).
\]
A lattice point with some $t(x_j)=0$ lies outside the closed support of $f_\iota$, so only sites with time coordinate at least $1$ contribute. Each such $\hat O_i(n)$ depends on edges in $\{t\ge0\}$. The series converges absolutely in the supremum norm, because the $\hat O_i$ are bounded and $f_\iota$ is a Schwartz function. So $F_\beta$ is bounded and depends only on edges in $\{t\ge0\}$. Since $\Theta\hat O_i(n)=\hat O_i(\theta n)$ and the $\hat O_i$ are real,
\[
  0\le\langle\overline{\Theta F_\beta}\,F_\beta\rangle_\beta=\sum_{\iota,\iota'}S^\beta_{(\iota',\iota)}\bigl((\Theta f_{\iota'})^*\otimes f_\iota\bigr)
\]
by H0. Each $(\Theta f_{\iota'})^*\otimes f_\iota$ lies in ${}^0\mathcal S$: inside each factor this holds by assumption, and at a point where a point $\theta$-reflected from the first group coincides with a point of the second group, the common time is either $\le0$ or $>0$, so one of the two factors vanishes on a neighbourhood. By Step~2 each term converges along $\beta_j$, and a finite sum of limits of non-negative numbers is non-negative. This gives E2, and positivity of the same form on families in $\mathcal T$.

\emph{Step 7: the gap inequality and cluster (E4).} For finite families $f,g$ in $\mathcal T$, with no constant component, put $S(f):=\sum_\iota S_\iota(f_\iota)$ and, for $s\ge0$,
\[
  B_{f,g}(s):=\sum_{\iota,\iota'}S_{(\iota',\iota)}\bigl((\Theta f_{\iota'})^*\otimes(g_\iota)_s\bigr)-\overline{S(f)}\,S(g),\qquad F_f(s):=B_{f,f}(s).
\]
By Step~6, $F_f(0)\ge0$ (apply E2 to the family consisting of $f$ and the constant $-S(f)$).

\begin{lemma}[gap inequality]\label{red:gap}
Assume H0, H1 with the unit $a=m$, H2 and H3. For finite families $f,g$ in $\mathcal T$ and $s\ge0$,
\[
  |B_{f,g}(s)|\le e^{-s}F_f(0)^{1/2}F_g(0)^{1/2};\qquad\text{in particular }|F_f(s)|\le e^{-s}F_f(0).
\]
\end{lemma}

\begin{proof}
On the lattice let $F^f_\beta,F^g_\beta$ be the observables of Step~6 and $v_f:=[F^f_\beta]-\langle F^f_\beta\rangle\Omega_\beta\in\Omega_\beta^\perp$, and similarly $v_g$. For $n\in\N$, $\tau_nF^g_\beta$ is the observable of the family $(g_\iota)_{na}$. Since $T_\beta$ is self-adjoint and $T_\beta\Omega_\beta=\Omega_\beta$, it maps $\Omega_\beta^\perp$ into itself, and
\[
  \langle v_f,T_\beta^nv_g\rangle=\langle\overline{\Theta F^f_\beta}\,\tau_nF^g_\beta\rangle_\beta-\overline{\langle F^f_\beta\rangle}\langle F^g_\beta\rangle=:B^\beta_{f,g}(na),
\]
the lattice-smeared analogue of $B_{f,g}$. Hence
\begin{equation}\label{red:latticegap}
  |B^\beta_{f,g}(na)|\le\|T_\beta|_{\Omega_\beta^\perp}\|^n\|v_f\|\|v_g\|=e^{-mn}\bigl(B^\beta_{f,f}(0)B^\beta_{g,g}(0)\bigr)^{1/2},
\end{equation}
and $mn=na$ because $a=m$. This uses only that $T_\beta$ is self-adjoint. Now let $n_j:=\lfloor s/a_j\rfloor$ and $s_j:=a_jn_j\to s$. Translation is continuous in each Schwartz norm, so $|(\Theta f_{\iota'})^*\otimes((g_\iota)_{s_j}-(g_\iota)_s)|_{N}\to0$, and the differences lie in ${}^0\mathcal S$. By the uniform bound \eqref{red:uniform} and Step~2, $B^{\beta_j}_{f,g}(s_j)\to B_{f,g}(s)$, and likewise at $s=0$. Pass to the limit in \eqref{red:latticegap}.
\end{proof}

The repaired step is this lemma. The source stated instead the monotone form $F(t+s)\le e^{-s}F(t)$ for all $t,s\ge0$. That form needs $T_\beta\ge0$, because for odd $n$ the number $\langle v,T_\beta^nv\rangle$ can be negative when $T_\beta$ has negative spectrum; H0 gives only $T_\beta$ self-adjoint with $T_\beta^2\ge0$. For the Wilson action $T_\beta\ge0$ is a theorem of \cite{Lu77}, which this paper does not use.

\emph{Cluster.} Let $f\in{}^0\mathcal S_+(\R^{4k})$, $g\in{}^0\mathcal S_+(\R^{4l})$ and let $\vec v\ne0$ be spatial. Translations by spatial vectors map ${}^0\mathcal S_+$ to itself and commute with $\Theta$. For $u,w$ in ${}^0\mathcal S_+$ write $\langle u,w\rangle:=S((\Theta u)^*\otimes w)$. By E2 this form is positive semidefinite, so it satisfies the Cauchy--Schwarz inequality, and $\langle w_{\lambda\vec v},w_{\lambda\vec v}\rangle=\langle w,w\rangle$ by E1.

(a) \emph{Reduction to compact support.} Let $\chi\in C_c^\infty(\R^{4k})$ equal $1$ near $0$ and put $f^R:=f\chi(\cdot/R)$. Then $f^R\to f$ in every Schwartz norm, $f^R$ has compact support contained in that of $f$, and $f^R\in{}^0\mathcal S_+$. Its support is a compact subset of the open set $\{0<t_1<\dots<t_k\}$, so it has positive distance from $\{t_1=0\}$ and from every coincidence set. By Cauchy--Schwarz, $|\langle f-f^R,g_{\lambda\vec v}\rangle|\le\langle f-f^R,f-f^R\rangle^{1/2}\langle g,g\rangle^{1/2}$, which tends to $0$ as $R\to\infty$ uniformly in $\lambda$, by E0$'$. The same holds for $g$, and $S(f^R)\to S(f)$ by E0$'$. So it suffices to treat $f,g$ with compact support of this kind.

(b) \emph{Rotation.} Let $\hat v:=\vec v/|\vec v|$ and let $\varrho\in SO(4)$ be the rotation by $\pi/2$ in the plane of $e_0$ and $\hat v$ that maps $\hat v$ to $e_0$. Put $u:=(\Theta f)^*\otimes g_{\lambda\vec v}$. By E1, $S(u)=S(u\circ\varrho^{-1})$, and $u\circ\varrho^{-1}=h\otimes(g')_{\lambda|\vec v|}$ with $h:=(\Theta f)^*\circ\varrho^{-1}$ and $g':=g\circ\varrho^{-1}$, both compactly supported, in a ball of radius $c_0$ say, and with supports away from all coincidence sets. Translating the whole configuration by $-c_0e_0$ (E1) gives $S(u)=S(h_{-c_0}\otimes(g'_{c_0})_{s})$ with $s:=\lambda|\vec v|-2c_0$. For $\lambda|\vec v|\ge2c_0$, $h_{-c_0}$ is supported in $\{$all $t<0\}$ and $G:=g'_{c_0}$ in $\{$all $t>0\}$. Put $\Phi:=(\Theta h_{-c_0})^*$. Then $\Phi,G\in\mathcal T$ and $h_{-c_0}=(\Theta\Phi)^*$.

(c) By Lemma~\ref{red:gap}, $S(u)=\overline{S(\Phi)}S(G)+B_{\Phi,G}(s)$ with $|B_{\Phi,G}(s)|\le e^{-s}F_\Phi(0)^{1/2}F_G(0)^{1/2}\to0$ as $\lambda\to\infty$. By E1 and reality, $\overline{S(\Phi)}=S((\Theta\Phi)^*)=S(h)=S((\Theta f)^*)$ and $S(G)=S(g)$. This is E4, with exponential rate at least $1$ for compactly supported $f,g$.

\emph{Step 8: not empty.} $S^\beta_{ii}(0,aN_\beta e_0)=1$ by the definition of $Z_i$, and $|aN_\beta-1|\le2a$ because $\xi-2<N_\beta\le\xi$ and $a\xi=1$. By \eqref{red:asym-eq} on $K_2$, $S_{ii}(0,e_0)=\lim_j\tilde S^{\beta_j}_{ii}(0,a_jN_{\beta_j}e_0)=1$. This is part~3. For later use, let $f\in C_c^\infty(\R^4)$ be real, non-negative, supported within distance $r_0<\tfrac14$ of $\tfrac12e_0$, with $\int f=1$. Then $f\in{}^0\mathcal S_+(\R^4)$, and as $r_0\to0$, $\langle f,f\rangle=\int\!\!\int S_{ii}(\theta x,y)f(x)f(y)\,dx\,dy\to S_{ii}(-\tfrac12e_0,\tfrac12e_0)=1$ by continuity and Step~3. The one-point functions vanish because the fields are centred, so $S_i(f)=0$.

\emph{Step 9: reconstruction and the gap.} By Steps~2--6, the $S_\iota$ satisfy E0$'$ and E1--E3, and Theorem~\ref{red:OS} gives a Wightman theory with fields $\varphi_1,\dots,\varphi_r$. In the description of the reconstructed space recalled after Theorem~\ref{red:OS}, a vector $\Psi=[f]$ with $f$ a finite family in ${}^0\mathcal S_+$ satisfies $\langle\Psi,e^{-sH}\Psi\rangle-|\langle\Omega,\Psi\rangle|^2=F_f(s)$. Let $\nu_\Psi$ be the spectral measure of $H$ for the component of $\Psi$ orthogonal to $\Omega$, so $F_f(s)=\int e^{-s\varepsilon}\,d\nu_\Psi(\varepsilon)$. If $\nu_\Psi([0,1-\eta])=c_1>0$ for some $\eta>0$, then $F_f(s)\ge c_1e^{-(1-\eta)s}$, which contradicts $F_f(s)\le e^{-s}F_f(0)$ (Lemma~\ref{red:gap}) for large $s$. So the spectral projection $E_H([0,1))$ annihilates the orthogonal component of a dense set of vectors, hence of every vector: $\operatorname{spec}(H|_{\Omega^\perp})\subset[1,\infty)$. By Step~8, the vector $[f]$ there is non-zero and orthogonal to $\Omega$, so $\Omega^\perp\ne\{0\}$.

\emph{Unique vacuum.} A translation-invariant vector is invariant under time translations, so $H\Psi=0$; by the gap, $\Psi$ is a multiple of $\Omega$.

\emph{Mass hyperboloid.} The Wightman theory carries a unitary representation of the Poincar\'e group, so the joint spectrum of $(H,\vec P)$ on $\Omega^\perp$ is a closed Lorentz-invariant subset of the closed forward cone. If it contained a point $p$ with $0<p^2<1$, a boost would carry $p$ to a point with energy $\sqrt{p^2}<1$. If it contained $p\ne0$ with $p^2=0$, boosts would give points of arbitrarily small energy. If it contained $p=0$, the energy $0$ would be in $\operatorname{spec}(H|_{\Omega^\perp})$. Each contradicts the gap. So the joint spectrum on $\Omega^\perp$ lies in $\{p^2\ge1,\ p^0>0\}$. This proves part~4 and completes the proof of Theorem~\ref{thm:conditional}. \qed

\emph{Where each hypothesis enters.} Step~1: H3(b), H3(c), H2. Step~2: H3(b), Step~1. Step~3: H0, H2, H3(c). Step~4: H0, H4, H3(c), Lemma~\ref{red:dense}. Step~5: reality and commutativity of the fields. Step~6: H0, site-scalarity, boundedness of the fields, Step~2. Step~7: H0 (self-adjointness of $T_\beta$), H1 with $a=m$, Steps~2 and~6, E1. Step~8: H3(a), H3(c), the normalization. Step~9: Theorem~\ref{red:OS} and the recalled description of the reconstructed space.

\subsection{What Theorem~\ref{thm:conditional} does not give}\label{red:not}

\begin{itemize}
\item \emph{Uniqueness.} The limit is taken along a subsequence. Different subsequences, or different choices of $\mu_\beta\in\cL_\beta$, need not give the same limit.
\item \emph{Interaction.} Nothing in H0--H4 excludes a free theory; Section~\ref{red:free} shows that the free lattice field satisfies the analogues of all of them. A test on composite fields does not detect interaction in the form ``the limit is not a generalized free field'': for the free field the Wick square $:\!\varphi^2\!:$ has non-zero truncated four-point function. A correct form would be a statement about the reconstructed theory, such as a non-trivial scattering matrix. It is not formulated here.
\item \emph{Short-distance behaviour.} The theorem says nothing about agreement with asymptotic freedom at short distances.
\item \emph{Other groups.} The proof uses no property of $\SU(2)$. For another compact simple group it requires H0--H4 for that group, and nothing is claimed about them here.
\end{itemize}
The official problem statement \cite{JW} remarks, in substance, that existence by compactness alone is not accepted unless other techniques establish the properties of the limit. Theorem~\ref{thm:conditional} is an argument of exactly that shape: Arzel\`a--Ascoli, limits of quadratic forms, and the reconstruction theorem. It assumes, as H1, H3 and H4, exactly the properties that such a remark asks to be established.

\subsection{A test of the chain on a free field}\label{red:free}

The lattice free massive scalar field on $\Z^4$, with mass parameter $m\to0$, gives an analogue of the setting of Theorem~\ref{thm:conditional}. The unit is the exact gap $a=E_0(m)=2\operatorname{asinh}(m/2)$, and the fields are $\varphi$ and $:\!\varphi^2\!:$. It is an analogue and not an instance: the fields are unbounded, there is no gauge group, and boundedness is used in Step~6. Moment bounds replace boundedness. We do not prove a version of Theorem~\ref{thm:conditional} for this analogue. The test checks that no hypothesis of the chain fails on a model whose answer is known, and that the chain then returns the free field of mass $1$, with two-point function $S_2(x)=K_1(|x|)/(|x|K_1(1))$.

Two separately written programs computed the numbers in Table~\ref{red:freetab}. Program~II uses a Bessel--Schwinger representation of the lattice Green function. Mass parameters are $m=\tfrac14,\tfrac18,\tfrac1{16},\tfrac1{32}$ unless a column says otherwise. These are numerical results, not proofs.

\begin{table}[ht]
\centering
\small
\begin{tabular}{@{}p{0.17\linewidth}p{0.47\linewidth}p{0.29\linewidth}@{}}
\toprule
Item & Program I & Program II \\
\midrule
H0 & RP matrix on 40 random points (seed 12345): least eigenvalue $+3.1\cdot10^{-10}$ & -- \\
H1 & zero-momentum correlator equals $e^{-E_0n}/(2\sinh E_0)$ to $10^{-13}$ ($n\ge1$), 4 of 4 masses & gap $2\operatorname{asinh}(m/2)$ to $10^{-15}$, 12 of 12 \\
H3(b), $\varphi$ & $\sup|x|^2S_2=1.64,\ 2.00,\ 2.11,\ 2.15$ & $1.636,\ 2.001,\ 2.114,\ 2.149$ \\
H3(b), $:\!\varphi^2\!:$ & -- & $\sup|x|^4S_{\varphi^2}\le4.62$ \\
H3(c) & Lipschitz constant on the axis, $|x|\ge\tfrac12$: $5.1,\ 10.0,\ 15.0,\ 19.1$ (continuum $25.1$) & -- \\
H4 & $|S(0,\cos1,\sin1,0)-S(0,1,0,0)|=3.9\cdot10^{-2},\ 2.7\cdot10^{-2},\ 6.7\cdot10^{-3}$; $|S(1,0,0,0)-S(\tfrac12,\tfrac12,\tfrac12,\tfrac12)|=3.7\cdot10^{-2},\ 9.0\cdot10^{-3},\ 2.2\cdot10^{-3}$ ($m=\tfrac18,\tfrac1{16},\tfrac1{32}$) & rotation differences agree with program I \\
Limit & relative error against $K_1(|x|)/(|x|K_1(1))$ at $m=\tfrac1{32}$: $\le2.4\cdot10^{-3}$ at 8 points, on and off axis & $\le2.3\cdot10^{-3}$ at 6 points; $a^{-2}G$ against $K_1(|x|)/(4\pi^2|x|)$: $\le1.4\cdot10^{-3}$ \\
$:\!\varphi^2\!:$, truncated 4-point & -- & Monte Carlo $1.54\cdot10^{-5}\pm2.1\cdot10^{-6}$; Wick $1.60\cdot10^{-5}$ \\
\bottomrule
\end{tabular}
\caption{Free-field analogue of the chain: numerical values.}\label{red:freetab}
\end{table}

\emph{Status of each line.} Wick's theorem for Gaussian measures is a theorem; with it, the truncated four-point function of $:\!\varphi^2\!:$ is $16$ times a sum of three cycle products of the positive two-point function, hence positive. The dispersion relation behind H1 and the gap $2\operatorname{asinh}(m/2)$ is a standard computation; here it is checked numerically. Reflection positivity of the nearest-neighbour Gaussian measure in site planes is standard and is not proved here; the eigenvalue in the table illustrates it. H3(b) for $\varphi$ follows from Wick's theorem, the inequality $|x_i-x_j|^{-2}\le\rho_i^{-1}\rho_j^{-1}$ and the bound $G_m(n)\le C/|n|^2$ on the lattice Green function, uniform in $m$; that bound is recalled from the literature on Green-function asymptotics, and we do not prove it here. Isotropy of the limit (H4) is supported by numerics only; its proof by a local central limit theorem is recalled and not written out. The convergence to $K_1(|x|)/(|x|K_1(1))$ is numerical. With these qualifications the free field satisfies the analogues of H0--H4, and the chain returns the free field of mass $1$. Interaction, in the sense of Section~\ref{red:not}, fails for $\varphi$ (its truncated four-point function vanishes).

\subsection{Status of the hypotheses for \texorpdfstring{$\SU(2)$}{SU(2)}}\label{red:statussec}

\begin{table}[ht]
\centering
\small
\begin{tabular}{@{}lll@{}}
\toprule
Hypothesis & Content & Status for $\SU(2)$ \\
\midrule
H0 (Hypothesis~\ref{hyp:H0}) & lattice axioms & proved \\
H1 = (A) (Hypothesis~\ref{hyp:A}) & gap at every coupling & open: it is the problem \\
H2 = (B) (Hypothesis~\ref{hyp:B}) & $m(\beta)\to0$ & proved (Theorem~\ref{thm:main}) \\
H3 (Hypothesis~\ref{hyp:H3}) & uniform regularity & open \\
H4 (Hypothesis~\ref{hyp:H4}) & one irrational rotation & open \\
\bottomrule
\end{tabular}
\caption{Status of the hypotheses of Theorem~\ref{thm:conditional} for $\SU(2)$.}\label{red:status}
\end{table}

For H0, translation and $W_4$ invariance of the periodic states pass to weak limits, and site-plane reflection positivity of the periodic Wilson states \cite{OS78} passes to weak limits. Hypothesis~(C) of Theorem~\ref{thm:reduction} is open for $\SU(2)$ as well.
